% Auto-generated by code/ — DO NOT EDIT BY HAND.
\begin{table}[htbp]
\centering
\caption{Robustness of the ASEAN Gini decline, 2014--2024.}
\label{tab:gini_robustness}
\begin{tabular}{lrrrrc}
\toprule
Specification & Regions & \shortstack{Gini\\2014} & \shortstack{Gini\\2024} & \shortstack{Change\\in Gini} & 95\% CI \\
\midrule
Headline (population-weighted, all units) & 278 & 0.584 & 0.528 & -0.057 & [-0.072, -0.042] \\
\midrule
Unweighted (equal weight per region) & 278 & 0.714 & 0.635 & -0.079 & [-0.089, -0.067] \\
Population fixed at 2014 & 278 & 0.584 & 0.528 & -0.057 & [-0.072, -0.042] \\
Excluding ever-dark units & 276 & 0.584 & 0.528 & -0.057 & [-0.072, -0.041] \\
Excluding Singapore & 269 & 0.565 & 0.516 & -0.050 & [-0.063, -0.036] \\
Excluding Myanmar & 261 & 0.572 & 0.512 & -0.060 & [-0.076, -0.044] \\
\bottomrule
\end{tabular}
\\[2pt]
\parbox[t]{\linewidth}{\footnotesize \emph{Notes:} All rows cover the ten ASEAN member states. The population surface in the source panel is a single epoch broadcast to every year, so weighting by each year's population and by 2014 population give the same series to within 3e-16; the row is reported to document that invariance, not as a distinct specification. Intervals from the paired region bootstrap, 2,000 draws.}
\end{table}
